\chapter{Ethical assessment of the project}\label{sec:ethics}

The development of autonomous driving systems cannot be separated from the ethical
weight of the domain in which they operate. A reinforcement learning agent that
controls a vehicle is, in a very literal sense, making decisions that affect human
safety, public infrastructure, and the welfare of other road users. This section
reflects on the ethical issues raised by the design of \textit{CALA\_SafeRL}, contrasts
the option of deploying a standard, unconstrained reinforcement learning agent against
the Shielded-RL approach actually adopted, and justifies the latter as the decision
that best satisfies the engineer's professional and ethical responsibilities. The
discussion closes by situating the project within the broader frameworks of digital
human rights and the United Nations Sustainable Development Goals (SDGs).

\section{Identification of ethical issues}\label{sec:ethics-issues}

Three ethical concerns are inherent to any project that trains a driving policy through
trial and error:

\begin{itemize}
  \item \textbf{Risk to human life and physical integrity.} Reinforcement learning is, by
        construction, an exploratory process: the agent improves by experiencing the
        consequences of its actions, including the negative ones. In a real or
        high-fidelity simulated traffic environment, ``negative consequences'' translate
        into collisions, off-road excursions, or dangerous manoeuvres near pedestrians and
        other vehicles. Even if the immediate setting is a simulator, the policies trained
        in it are intended to generalise toward real-world deployment, so any tolerance for
        unsafe exploratory behaviour during training is, in principle, a tolerance for
        unsafe behaviour that could eventually reach public roads.

  \item \textbf{Algorithmic accountability.} A policy represented as a neural network is
        difficult to audit. When such a system causes harm, it is not always possible to
        trace the decision back to an interpretable cause, which complicates the
        attribution of responsibility among the developer, the operator, and the
        manufacturer. An engineering team that deploys a purely learned controller without
        any verifiable safety layer is, in effect, delegating safety-critical decisions to
        a component whose failure modes cannot be fully characterised in advance.

  \item \textbf{Black-box decision making.} Beyond accountability after the fact, there is
        a forward-looking concern: end users, regulators, and the general public are
        increasingly reluctant to trust systems whose behaviour cannot be bounded or
        explained. A driving agent that "mostly works" but offers no guarantee about its
        worst-case behaviour is difficult to certify and difficult to trust, regardless of
        how well it performs on average.
\end{itemize}

These three issues are not abstract concerns external to the project; they are direct
consequences of the choice of algorithm (PPO, a model-free reinforcement learning
method) and of the choice of action interface (continuous steering and throttle/brake
commands sent directly to the simulated vehicle).

\section{Critical analysis of design options}\label{sec:ethics-options}

Two broad design philosophies were considered during the early stages of the project,
and the trade-off between them is fundamentally an ethical one rather than a purely
technical one.

\paragraph{Option A: an unconstrained reinforcement learning agent.}
The first option is to train the PPO agent directly against the CARLA simulator, with no
intermediate safety layer (the \texttt{--shield\_type none} configuration evaluated in
this work, Section~\ref{ch:experimentation}). This option has clear technical
advantages: the agent's gradient signal is unfiltered, exploration is unconstrained, and
in principle the policy can discover driving strategies that a hand-designed corrector
might not anticipate. From a purely performance-driven perspective, this option could be
argued to maximise \textit{learning speed} and \textit{asymptotic performance}.

However, this option places the entire burden of safety on the learning process itself,
during a phase in which the policy is, by definition, not yet competent. Early in
training the agent behaves close to randomly; under Option A, that randomness is applied
directly to the vehicle's controls, and the resulting collisions and off-road events are
treated merely as numerical penalties in a reward function. There is no mechanism that
prevents an unsafe action from being executed --- safety is, at best, an emergent
property that the agent might converge towards, not a property that is guaranteed at any
point during training or deployment. The evaluation data gathered for this project
(Chapter~\ref{ch:experimentation}) confirms that, even after substantial training, a
non-trivial fraction of episodes under the unshielded configuration end in a crash or an
off-road event, which would be unacceptable if those episodes corresponded to interactions
with real vehicles, pedestrians, or infrastructure.

\paragraph{Option B: Shielded-RL with a deterministic safety layer.}
The second option, and the one adopted in this project, interposes a \textit{safety
shield} between the agent and the simulator. The shield --- in its most developed form,
the \texttt{CarlaAdaptiveHorizonShield} described in Chapter~\ref{ch:design} ---
monitors the agent's proposed action, predicts the vehicle's short-horizon trajectory
using a kinematic bicycle model, and deterministically overrides the action whenever the
prediction indicates an unacceptable risk (collision, lane departure, or deadlock). The
agent still learns from experience, including from the cases in which its action was
corrected, but the action that actually reaches the simulator is bounded by an explicit,
inspectable, rule-based safety envelope.

This option does not maximise raw learning speed or asymptotic reward in the same
unconstrained sense as Option A: the shield occasionally prevents the agent from taking
the action it would have chosen, and the credit-assignment machinery
(Chapter~\ref{ch:design}) must be designed carefully so that the agent is not
penalised for the shield's intervention. What Option B does provide, by construction, is
a \textit{guaranteed upper bound on the severity of unsafe behaviour} at every stage of
training and deployment, together with a fully interpretable component (the shield) whose
decisions can be logged, audited, and explained independently of the neural network.

\paragraph{The ethical trade-off.}
Stated in its simplest form, the choice between Option A and Option B is a choice between
\textit{optimising for performance under an implicit assumption that safety will emerge},
and \textit{optimising for performance under an explicit constraint that safety is
guaranteed by design}. The first option treats safety as a metric to be learned; the
second treats safety as a precondition that is never violated, regardless of what the
learned component has or has not yet learned. From an engineering ethics standpoint,
these are not equivalent options separated by a small performance gap --- they represent
two fundamentally different attitudes toward the question of who, or what, is responsible
for preventing harm.

\section{Justification of the selected decision}\label{sec:ethics-justification}

The Shielded-RL architecture (Option B) was selected, and this decision is justified
along the three axes specified for this assessment.

\paragraph{Engineering professional ethics.}
Professional codes of conduct for engineers --- including those referenced by the
\textit{Colegio de Ingenieros} bodies and by international bodies such as the IEEE and
ACM --- consistently place the protection of public health, safety, and welfare above
all other professional obligations, including those toward employers, clients, or the
advancement of the technology itself. In the context of this project, that principle
translates directly into a design requirement: \textit{no action that the simulated
vehicle executes should be allowed to cause a collision or an off-road excursion if a
safe alternative exists and can be computed in real time}. The
\texttt{CarlaAdaptiveHorizonShield} operationalises this requirement. Its risk
classification (safe / warning / critical, Chapter~\ref{ch:design}) and its
trajectory-prediction logic exist precisely so that the "first do no harm" principle is
enforced at the level of the control loop, not merely hoped for as a statistical outcome
of training. Choosing to build and integrate this component, rather than relying on the
learned policy alone, is a direct application of the duty to prioritise public safety
over project convenience or raw performance metrics.

\paragraph{Social impact and the role of deterministic guardrails in building trust.}
Public acceptance of autonomous systems --- whether autonomous vehicles or, more broadly,
AI systems embedded in everyday infrastructure --- depends heavily on the existence of
\textit{verifiable, deterministic guardrails} that bound the system's worst-case
behaviour. This is the same principle that motivates safety filters and content
guardrails around large language models: a system that is occasionally impressive but
unboundedly unpredictable does not earn public trust, whereas a system whose
worst-case behaviour is provably constrained, even if its average behaviour is more
modest, can be reasoned about, certified, and ultimately trusted. The shield architecture
in this project plays exactly this role for the driving agent. Because the shield's
corrective logic is rule-based and inspectable (kinematic predictions, explicit distance
and time-to-collision thresholds), its behaviour can be explained to a non-specialist
in terms of physical quantities (``the vehicle was predicted to leave its lane within
\(N\) steps, so its steering was corrected''), rather than in terms of opaque network
activations. This interpretability is, in itself, a contribution to the social
acceptability of learned autonomous controllers: it demonstrates that performance and
verifiable safety are not mutually exclusive design goals, and that a learning system
can be embedded inside --- rather than substituted for --- a deterministic safety
framework.

\paragraph{Alignment with Deusto's values and the Sustainable Development Goals.}
This project is also consistent with the broader institutional commitments of the
University of Deusto, in particular the principles articulated in the
\textit{Declaración Deusto de Derechos Humanos en Entornos Digitales}, which calls for
the design of digital and automated systems that actively safeguard human rights ---
including the right to life and physical integrity --- rather than treating their
protection as an afterthought. A driving system whose safety properties are designed in
from the outset, and whose decisions remain auditable, is a direct application of that
principle to the domain of autonomous mobility.

The project's goals also align with several of the United Nations Sustainable
Development Goals:

\begin{itemize}
  \item \textbf{SDG 3 (Good Health and Well-being):} road traffic collisions are a
        leading cause of injury and death worldwide. By developing and evaluating
        mechanisms that reduce the rate of crashes and off-road events during the
        learning process of an autonomous agent (Chapter~\ref{ch:experimentation}), this
        project contributes, even at a research scale, to the broader goal of reducing
        traffic-related harm.

  \item \textbf{SDG 9 (Industry, Innovation, and Infrastructure):} the project advances
        the state of the art in safe reinforcement learning for autonomous driving,
        proposing and evaluating a concrete architecture (the adaptive horizon shield)
        that other researchers and practitioners can build upon, fostering responsible
        innovation in intelligent transport infrastructure.

  \item \textbf{SDG 11 (Sustainable Cities and Communities):} autonomous vehicles are
        widely expected to play a role in future urban mobility systems. Safety
        mechanisms such as the one developed in this project are a prerequisite for the
        responsible integration of such vehicles into shared urban spaces, where
        pedestrians, cyclists, and other vehicles must coexist safely with autonomous
        agents.
\end{itemize}

\paragraph{Summary.}
Taken together, these considerations support the conclusion that the Shielded-RL
architecture is not merely a technical refinement of the baseline PPO agent, but the
ethically appropriate design choice for this project. It reflects the professional duty
to prioritise public safety over unconstrained optimisation, it embodies the kind of
deterministic, interpretable guardrail that is increasingly recognised as a precondition
for public trust in autonomous systems, and it situates the project within the human
rights and sustainability commitments that frame the work of engineers trained at the
University of Deusto.
